\chapter{Что нужно построить}

\begin{warnbox}[Глава заменена роадмапом экосистемы]
Состояние на 29 сентября 2026 года. Что строится и в каком порядке, теперь задают роадмап экосистемы (вики, \texttt{roadmap/ecosystem-roadmap.ru.md}, черновик v4: этапы M0--M8, треки PC, KP и CM) и спецификации модулей: конфигурация процесса (curator-playbook, трек PC), хост сессий (agent-session-host, этап SH), главный цикл узла (curator-keeper, трек KP), доверие (curator-trust), подписанные сообщения и аренды (waggle, трек CM). Распределённый пул идёт треком AR в curator-agent-runtime после этапов M4 и M7. Демон среды, секвенсер журнала и вход в санбокс пока описаны только в работе. Минимальный runtime отложен вместе с собственным агентским циклом (curator-keeper, §7.7). Текст главы сохранён как история; где он расходится с роадмапом или спецификацией модуля, правы они.
\end{warnbox}

Три новых компонента и три расширения существующих. Для каждого: ответственность, интерфейс в набросках и то, какие инварианты он обеспечивает. Наброски интерфейсов здесь не нормативны; нормативные документы со схемами и векторами соответствия пишутся уровнем ниже.

\section{Демон среды}

Демон среды есть провайдер исполнения на хосте: единственный процесс на среде, который принимает команды извне. Он принадлежит проекту, его хеш опубликован в манифесте, он не слушает портов и работает по стандартным потокам поверх SSH.

Ответственность: проверка ключа ноды и записи назначения; CA среды и выпуск сертификатов на задачу с принудительной командой; приём соединений (OpenSSH с сертификатами на L0, встроенный SSH-сервер в самом демоне для публичного роя); создание и уничтожение санбоксов; исполнение команд в санбоксе с лимитами и receipts; восстановленные сессии; фоновые процессы и spill; построение интеграционного кандидата; атомарное продвижение головы; операции раннера; светлый клиент журнала с защитой от отката; безопасный режим по сертификату остановки; инвентарь ограничений.

\begin{exbox}[{Поверхность демона, набросок}]
\begin{lstlisting}
permit.verify(record, proof)       -> ok | reject(reason)
cert.issue(assignment)             -> ssh_cert(principal=task, ttl=window, force_command)
sandbox.open(task, plan_digest)    -> sandbox_id
sandbox.close(sandbox_id)
exec(sandbox_id, session, cmd, limits) -> receipt
bg.start / bg.status / bg.log / bg.kill
files.read_range / files.apply_patch
candidate.build(task, R)           -> {M, I}
main.cas(M, I, integration_record) -> ok | stale(M_now)
runner.call(op_id, idem_key, op, args, record) -> outcome(done|rejected|unknown)
runner.reconcile(op_id)            -> outcome
stop.apply(certificate)            -> safe_mode
inventory()                        -> {fs: full, net: full, mem: partial, ...}
\end{lstlisting}
\end{exbox}

Обеспечивает инварианты I1 (включая проверку актуальности на границе исполнения), I3, I6 в части среды, I8 в части высшей версии, I9, I10.

Платформенные бэкенды: macOS через Seatbelt и клоны APFS с группой процессов и сторожем; Linux через bubblewrap с cgroups, с возможностью gVisor или Firecracker для чужих агентов. Бэкенд выбирается конфигурацией среды и отражается в инвентаре.

Планировщик слотов: очередь назначений, слоты как ёмкость, приоритет по эпохе и оценке. Egress-прокси: локальный HTTP-прокси с allowlist из манифеста, единственный сетевой выход из санбокса.

\begin{warnbox}[{Четыре места, где легко ошибиться}]
Первое: дать агенту любой путь в shell хоста, хоть на секунду, хоть для отладки. Второе: построить кандидата, а одобрение или проверки взять от исходной ревизии. Третье: принять снапшот журнала без проверки высшей версии. Четвёртое: повторить операцию раннера с исходом не установлен без сверки. Каждая ошибка обнуляет соответствующий инвариант, и первые три из них исполняемая модель (раздел 11.8) находит как нарушения.
\end{warnbox}

\section{Секвенсер журнала}

Секвенсер расширяет реестр аудита до журнала. Ответственность: приём подписанных записей от оркестратора, воркеров и кор-нод; присвоение порядка; проверка правил валидаторов на уровне сервиса (для ускорения, но не вместо подписей); сбор эндорсментов кор-нод до порога кворума; выпуск снапшотов; доказательства включения; сертификаты доступности; сертификат остановки как запись с особым правилом.

\begin{exbox}[{Расширение реестра, что добавляется}]
\begin{itemize}
\item схемы записей по типам из главы 5 вместо записи аудита артефакта, включая версии контракта, зависимости, кандидата и адресованные к нему одобрения и аттестации с тремя исходами;
\item поле кворума: список эндорсментов кор-ключей, запись совершена при достижении порога уровня;
\item метод \texttt{inclusion\_proof(entry)} и снапшот с корнем для светлых клиентов;
\item метод \texttt{availability(hash)}: реплики подтверждают хранение, сертификат прикладывается к записи;
\item тип записи остановки, принимаемый с одной подписью кор-ключа, и возобновления, принимаемого с кворумом;
\item на L2: интерфейс к движку консенсуса вместо собственного порядка.
\end{itemize}
\end{exbox}

Обеспечивает I4 (правила исполняются как детерминированные предикаты), I7 (эпоха и токен в записях), I8 (остановка и версии), и хранит материал для I2, I3, I5.

\section{Протокол ноды}

Протокол ноды есть единственный способ для контрибьютора говорить с проектом, одинаковый для обеих форм клиента.

\begin{exbox}[Сообщения]
\begin{lstlisting}
join(manifest_url, node_key, axes, identity_proofs) -> plan, authority_diff
heartbeat(capacity_left, reset_window, cost_class, slots_wanted)
assign(task, role, plan_digest, sandbox_endpoint)   <- from project
session.open / exec / close                 -> via environment daemon
submit(task, R, B, receipts_head, signature)
review.request(task, I_or_R)                <- from project
review.submit(task, verdict, findings_hash, signature)
attest.request(task, I) / attest.submit(task, I, result, signature)
leave()
\end{lstlisting}
\end{exbox}

Присоединение компилирует план из манифеста проекта и показывает разницу полномочий. Heartbeat несёт остаток ёмкости, окно сброса и класс стоимости. Подписи ставятся брокером на машине ноды.

Обеспечивает I5 в части объявления осей, I6, и даёт секвенсеру данные для I2.

\section{Резолверы идентичности}

Слой идентичности (глава 8) в MVP не входит, но интерфейс стоит зафиксировать сейчас, чтобы протокол ноды и типы записей его не исключали. Резолвер есть плагин с одним методом: по строке идентичности вернуть набор отпечатков ключей и свидетельство. Первые три: TXT-запись домена под DNSSEC, well-known по HTTPS, ключи аккаунта на хостинге кода.

\begin{exbox}[Интерфейс резолвера]
\begin{lstlisting}
resolve(identity) -> {fingerprints: [..], evidence_hash, fetched_at, resolver_id}
                   | unavailable(reason)
\end{lstlisting}
Запись подтверждения идентичности содержит идентичность, класс доказательства, хеш свидетельства и окно; она совершена, когда $t$ валидаторов, каждый резолвивший сам, записали одинаковый хеш. Запись сертификации ключа содержит ключ ноды, подписанный ключом идентичности, с началом и концом действия. Файл допустимых подписантов git для среды выводится из этих двух типов записей.
\end{exbox}

\section{Скилл удалённого shell}

Скилл есть пакет Curator: инструкции для модели (минимальные: у вас есть shell в рабочем окружении), провайдер, который говорит с демоном среды, и артефакты. Провайдер реализует восстановленную сессию (снимок состояния по отдельному дескриптору, восстановление перед командой), именованные сессии с сериализацией мутирующих вызовов, фоновые процессы по дескриптору, spill больших выводов, отмену по группе процессов, переподключение без потери состояния сессии.

Отличие от демо на одноразовой машине: демон принадлежит проекту, не заливается контрибьютором, а пинится по хешу из манифеста. Провайдер проверяет, что подключился к тому демону, идентичность которого объявлена, и это доверие к проекту, а не верификация.

Вторая привязка публичного клиента, подпись ревизий, есть операция брокера: подписать хеш, который совпадает с хешем в журнале задачи.

\section{Минимальный runtime}

Публичный клиент есть runtime с планом из двух привязок и командой присоединения. Установка одной командой: ставится runtime, по адресу манифеста компилируется план, закрывается closure с пинами скилла shell, скилла подписи и хеша демона, показывается разница полномочий, и это вся установка. Работа над тем, чтобы весь набор ставился одной командой и был бесшовным для контрибьютора, уже идёт.

\section{Оркестратор поверх спавна борда}

Оркестратор есть кор-нода в роли: агент, который читает цели, пишет задачи в борд, отправляет записи в секвенсер и назначает ноды. Спавн борда уже умеет запускать роли разработчика и ревьювера в копиях дерева с командами проверки; оркестратор заменяет локальный запуск назначением на ноду через протокол, добавляет планировщик по ёмкости и областям, заморозку кандидатов и случайный выбор ревьювера, меню при превышении $N$, аренду с heartbeat и чекпоинты.

\section{Два эксперимента с первого дня}

Соблюдение протокола и полезность роя есть разные вопросы, и они проверяются по-разному.

\term{Соблюдение протокола.} Реализация подвергается сбоям и нарушителям, воспроизводимым программно, без участия моделей: рабочая нода, выходящая за write scope; запись оркестратора с устаревшей арендой; подмена кандидата после аттестации; изменение критериев после ревью; повторная доставка записи об интеграции; операция раннера с потерянным исходом; смерть оркестратора с удалением его сессии. Критерий: ни одно нарушение не приводит к эффекту на среде или продвижению головы, а восстановление проходит по журналу. Последний сценарий, остановить оркестратора, удалить его сессию и проверить, продолжит ли преемник из журнала и доступных артефактов, есть прямая проверка инварианта I7 и утверждения о достаточности внешнего состояния.

\term{Полезность роя.} Единица сравнения есть заранее определённая внешняя задача или завершённая функция проекта с одинаковыми критериями приёмки, а не внутренние задачи оркестратора: он сам декомпозирует работу и может менять их число, не меняя объёма результата. Три условия при сравнимых ресурсных ограничениях: одиночный агент; параллельный процесс без протокола, агенты одного владельца с общим репозиторием и блокировками; развёртывание L0. Стоимость включает inference, исполнение проверок, занятость среды и человеческое вмешательство. Диагностические метрики собираются из журнала и receipts без дополнительной инструментовки: принятые задачи на миллион токенов, время от авторизации до интеграции, ожидание слотов, блокировки областей, повторная работа после смены головы, восстановление после потери агента, доля пересмотренных отклонений, доля откатов через $k$ эпох, часы человеческого ревью на принятую задачу.

Дашборд поверх журнала есть отдельная маленькая задача; важно, чтобы данные для него существовали с первого совершённого назначения. Ориентир масштаба для первого числа: эксперимент, в котором три личные подписки за три дня совместно формализовали крупный результат.

\section{Исполняемая модель фрагмента}

Инварианты и правила валидаторов детерминированы, и ключевой фрагмент протокола допускает исполняемую модель. К третьей редакции работы приложена модель фрагмента кандидат, аттестации, запись об интеграции, атомарное продвижение головы, проверенная исчерпывающим перебором состояний в заданных границах.

\begin{exbox}[Что в модели]
Две задачи: $A$ с полным поведением (контракт может менять версию, записи могут нести устаревшую эпоху) и $B$, которая только интегрируется и тем сдвигает голову. Состояние: голова, эпоха аренды и для каждой задачи наличие ревизии, база кандидата, аттестация (адрес и версия контракта), версия контракта, признак интеграции, запись об интеграции с числом доставок. Действия: представить ревизию; построить кандидата над текущей головой; аттестовать кандидата; изменить версию контракта; сменить эпоху; предложить запись с текущей или устаревшей эпохой; доставить запись, в том числе повторно. Границы: не больше двух версий контракта, двух эпох, двух доставок на запись. Правила при доставке: голова равна базе записи; запись несёт текущую эпоху; аттестация ссылается на текущую версию контракта; аттестация адресована ровно к кандидату. Свойство: голова продвигается к кандидату, построенному над головой в этот момент, с аттестацией к этому кандидату и текущему контракту, по записи текущей эпохи, и ни одна запись не применяется дважды.
\end{exbox}

При всех четырёх правилах достижимо 1090 состояний, свойство выполнено во всех. При отключении любого одного правила перебор находит контрпример: без проверки головы одна запись применяется дважды при повторной доставке, а также голова продвигается к кандидату, построенному не над текущей головой, после интеграции другой задачи; без проверки эпохи запись устаревшей аренды продвигает голову; без проверки версии контракта аттестация под прежними критериями принимается после их изменения; без привязки к кандидату аттестация, сделанная на базе до сдвига головы, переносится на кандидата над новой головой. Последний контрпример есть ровно сценарий переноса одобрения по patch-id, и именно он вывел это правило из гарантии.

\begin{warnbox}[Границы результата]
Это проверка ограниченного пространства сценариев, а не доказательство для всех допустимых исполнений: границы заданы явно, среда и валидаторы в модели честны, суждение и содержимое кода не моделируются. Что модель даёт реализации: её трассы становятся векторами соответствия демона среды и секвенсера, и первый набор векторов получается запуском модели, а не ручным написанием. Следующий шаг есть та же модель в языке спецификаций с механической проверкой и распространение на аренду и остановку.
\end{warnbox}
